% ============================================================================
% Reporte de trabajo en equipo - Entrega 1
% Documento SEPARADO del reporte principal, con limite de 1 pagina, segun la
% seccion "Entregables" de maia_pds_proy_s2.pdf.
% ============================================================================
\documentclass[11pt]{article}

\usepackage[utf8]{inputenc}
\usepackage[T1]{fontenc}
\usepackage[spanish,es-tabla]{babel}
\usepackage{geometry}
\geometry{a4paper, top=2.2cm, bottom=2.0cm, left=2.4cm, right=2.4cm}

\usepackage{titlesec}
\usepackage{booktabs}
\usepackage{array}
\usepackage[table]{xcolor}
\usepackage{hyperref}
\usepackage{parskip}

% Misma paleta que el reporte principal, para que ambos documentos se lean
% como parte de una sola entrega.
\definecolor{primaryblue}{RGB}{28,58,92}
\definecolor{accentblue}{RGB}{47,109,161}
\definecolor{midgray}{RGB}{110,110,110}

\titleformat{\section}
  {\large\bfseries\color{primaryblue}}
  {\thesection.}{0.4em}{}
\titlespacing*{\section}{0pt}{1.0em}{0.4em}

\hypersetup{
  colorlinks=true,
  urlcolor=accentblue,
  pdftitle={Reporte de trabajo en equipo - Entrega 1},
}

\pagestyle{empty}  % una sola pagina: el numero de pagina sobra

\begin{document}

% ----------------------------------------------------------- encabezado
\begin{center}
{\large\bfseries\color{primaryblue} Reporte de trabajo en equipo}\\[0.15cm]
{\color{midgray} Recomendación local de citas académicas mediante aprendizaje automático}\\[0.15cm]
{\color{midgray} Microproyecto -- Entrega 1 · Proyecto: Desarrollo de soluciones}\\[0.15cm]
{\color{midgray} Maestría en Inteligencia Artificial -- Universidad de los Andes}\\[0.15cm]
{\color{midgray} 23 de agosto de 2026}
\end{center}

\vspace{0.2cm}
{\color{primaryblue}\rule{\linewidth}{0.8pt}}
\vspace{0.2cm}

% ----------------------------------------------------------- organizacion
\section{Organización del trabajo}

El equipo distribuyó las actividades de la Entrega 1 entre la configuración de
la infraestructura reproducible, la exploración de los datos, la documentación y
la adecuación del reporte. La coordinación se realizó de forma asíncrona,
integrando cada aporte al repositorio mediante ramas de trabajo y \textit{Pull
Requests} revisados antes de incorporarse a la rama principal.

Las contribuciones fueron verificadas mediante el historial de \textit{commits}
del repositorio Git, de acuerdo con el requisito de evidenciar el aporte
individual de cada integrante:

\begin{center}
\url{https://github.com/Byron971/microproyecto-local-citation}
\end{center}

% ----------------------------------------------------------- aportes
\section{Aportes individuales}

\begin{table}[h!]
\centering
\small
\begin{tabular}{p{3.0cm} p{11.5cm}}
\toprule
\textbf{Integrante} & \textbf{Evidencia registrada} \\
\midrule
John Arias &
Configuración inicial de DVC y estructura base; versionamiento del
\textit{dataset}; remoto AWS S3; \textit{notebook} inicial de EDA; análisis de
relaciones del \textit{dataset} e integración de contribuciones mediante
\textit{Pull Request}. \\[3pt]

José Zambrana &
Configuración de remoto DVC; actualización amplia del \textit{notebook} de EDA;
incorporación de gráficas y evidencias; dependencias; creación y reorganización
del reporte para alinearlo con la rúbrica. \\[3pt]

Carlos Caicedo &
Mejoras del informe y de las evidencias DVC/S3; eliminación de contenido
redundante; corrección de la maqueta; recuperación y ajuste de la arquitectura
conceptual y del resultado esperado. \\[3pt]

Carlos Duque &
Actualización de la documentación del proyecto mediante el \texttt{README} e
incorporación de archivos asociados a la definición y reproducibilidad del
entorno, incluyendo \texttt{.python-version}, \texttt{pyproject.toml} y
\texttt{uv.lock}. \\
\bottomrule
\end{tabular}
\end{table}

% ----------------------------------------------------------- verificacion
\section{Verificación de la reproducibilidad}

El historial del repositorio evidencia contribuciones de los cuatro integrantes.
Adicionalmente, durante la validación final se comprobó la reproducibilidad del
proyecto mediante la creación de un entorno limpio, la recuperación de los datos
versionados desde Amazon S3 mediante DVC y la ejecución completa del
\textit{notebook} de exploración sin errores.

El plan de trabajo de las Entregas 2 y 3 quedó registrado en el repositorio como
tareas individuales agrupadas en dos hitos con sus fechas límite, cada una con su
criterio de aceptación, de modo que la asignación y el avance de cada integrante
puedan verificarse directamente sobre el repositorio.

\end{document}
